% Generated draft: requires mathematical and source review.
\documentclass[11pt]{article}
\usepackage[margin=25mm]{geometry}
\usepackage{fontspec}
\setmainfont{Latin Modern Roman}
\usepackage{amsmath,amssymb,mathtools,unicode-math}
\setmathfont{Latin Modern Math}
\usepackage{enumitem,needspace,xurl,hyperref}
\setlength{\parindent}{0pt}
\setlength{\parskip}{5pt}
\setlength{\emergencystretch}{3em}
\newcommand{\N}{\mathbb{N}}
\newcommand{\Z}{\mathbb{Z}}
\newcommand{\Q}{\mathbb{Q}}
\newcommand{\R}{\mathbb{R}}
\newcommand{\C}{\mathbb{C}}
\newcommand{\F}{\mathbb{F}}
\newcommand{\A}{\mathbb{A}}
\newcommand{\PP}{\mathbb{P}}
\newcommand{\E}{\mathbb{E}}
\newcommand{\eps}{\varepsilon}
\newcommand{\dd}{\,\mathrm{d}}
\begin{document}
\section*{UnsolvedMath ID: 30006334}
\textbf{Source problem code:} catalogue entry
\subsection{Definitions and mathematical statement}
The supplied record attributes the following conjecture to Nekov\'a\v{r}--Scholl. Write
\[
C=R\Gamma(\operatorname{Sh}_\infty,L_\xi)
   \in D^b(\Gamma_E,\Q_\ell).
\]
Here \(R\Gamma\) denotes derived global sections, and \(D^b\) denotes a bounded derived category. The source uses \(\operatorname{Sh}_\infty\) for a Shimura-variety tower and \(L_\xi\) for its coefficient system. It does not supply the underlying Shimura datum, the definition of the tower or coefficient system, the fields \(E,F\), the precise plectic group \(\Gamma_{F/\Q}\), or the representation-category conventions defining these derived categories. These data must be supplied before the conjecture becomes a fully specified mathematical assertion.

The requested lifting statement can be expressed conditionally as follows. Let
\[
\operatorname{Res}_{[\mu]}:
D^b(\Gamma_{F/\Q},\Q_\ell)
\longrightarrow D^b(\Gamma_E,\Q_\ell)
\]
stand for the restriction functor intended by the source's phrase ``via \([\mu]\).'' Its defining group homomorphism and dependence on \([\mu]\) are not provided in the record. For every instance covered by the intended, but unspecified, hypotheses, the conjecture asks for a canonical object
\[
\widetilde C\in D^b(\Gamma_{F/\Q},\Q_\ell)
\]
and a canonical isomorphism
\[
\alpha:\operatorname{Res}_{[\mu]}\widetilde C
\xrightarrow{\ \sim\ } C,
\]
compatible with the Hecke actions. In particular, for each intended Hecke operator \(T\), its lifted action \(\widetilde T\) should satisfy
\[
\alpha\circ\operatorname{Res}_{[\mu]}(\widetilde T)
=T\circ\alpha.
\]
The Hecke algebra, its allowed operators, and the precise meaning of canonicity are not specified. This conditional expression records the requested lift; it does not supply the missing hypotheses or a uniqueness criterion.

The source partly motivates the conjecture by applications to semisimplicity of Shimura-variety cohomology as Galois representations, citing its references [4] and [5]. It also describes an \emph{in-progress} result [12]: for distinct primes \(p\ne\ell\) and Shimura data of abelian type, a canonical action of a local plectic reflex Galois group, more precisely a Weil-group version, is claimed on the corresponding cohomology. This is a qualified local claim, not a supplied proof of the displayed global conjecture. For its explanatory formalism, the excerpt fixes a torsion \(\Z_\ell\)-algebra \(\Lambda\), uses constant coefficients, and says that \(\Q_\ell\)-coefficients can be obtained by an inverse limit and base change. It discusses the \(v\)-stack \(\operatorname{Bun}_G\) over \(\mathbf F_p\) of bundles on the Fargues--Fontaine curve and its derived category of \(\Lambda\)-valued \(\acute{e}\)tale sheaves. It attributes Hecke endofunctors with internal \(W_{\Q_p}\)-actions, or actions of a subgroup, to Fargues--Scholze [2], and a cohomology formula using abelian-type Igusa stacks to [3]. The excerpt ends during another in-progress theorem, so its full statement is unavailable.
\subsection{Short English statement}
The Nekov\'a\v{r}--Scholl conjecture asks whether the derived cohomology of a Shimura-variety tower has a canonical plectic Galois action extending its given Galois action via \([\mu]\), compatibly with Hecke operators. The supplied excerpt omits the setup needed to specify this extension precisely. Its local abelian-type result is explicitly described as work in progress.
\subsection{Sources}
\begin{itemize}
\item Local UnsolvedMath registry: supplied source record, permanent ID \(30006334\), exact source problem code \texttt{catalogue entry}, title ``The Nekov\'a\v{r}--Scholl Plectic Cohomology Conjecture.'' Attribution is to Nekov\'a\v{r}--Scholl in the excerpt; the registry's proposer field is unspecified.
\item \emph{Non-Archimedean Geometry and Applications} (2025), Oberwolfach Reports, pp.~1517--1578, DOI \texttt{10.4171/owr/2025/29}. Supplied report provenance and literature-source link: \url{https://doi.org/10.4171/owr/2025/29}.
\item \emph{The plectic conjecture over local fields}. Supplied literature-source link: \url{https://arxiv.org/abs/2208.09962}. The record supplies no authors or further bibliographic details for this reference.
\item The literature assessment dated \texttt{2026-08-22} is supplied dated curation, not independently verified here. It labels the conjecture open and describes it as widely open, with results for CM points and partial cases but no general canonical lift compatible with Hecke actions. The background explicitly requests independent verification of time-sensitive status and cited claims.
\item The registry labels the statement \texttt{corrected\_verified}, explaining that only formula delimiters, notation, grammar, and typography were repaired. This records the supplied editorial qualification and does not establish mathematical correctness. The excerpt's numbered references [2], [3], [4], [5], and [12] have no accompanying bibliography in the supplied record.
\end{itemize}
\end{document}
